\documentclass[12pt,a4paper]{article}
\usepackage[T2A]{fontenc}
\usepackage[utf8]{inputenc}
\usepackage[bulgarian]{babel}
\usepackage{lmodern}
\usepackage{amsmath}
\usepackage{amssymb}
\usepackage{mathtools}
\usepackage[margin=2.2cm]{geometry}
\usepackage{titlesec}
\usepackage{enumitem}
\usepackage{array}
\usepackage{booktabs}
\usepackage{longtable}
\usepackage[table]{xcolor}
\usepackage{tikz}
\usepackage{fancyhdr}
\usepackage[hidelinks]{hyperref}
\usepackage{tabularx}

\definecolor{titleblue}{RGB}{20,40,90}
\definecolor{sectionblue}{RGB}{30,60,120}
\definecolor{rulegray}{RGB}{150,150,150}

\titleformat{\section}
  {\large\bfseries\color{sectionblue}}{\thesection}{1em}{}
\titleformat{\subsection}
  {\normalsize\bfseries\color{sectionblue}}{\thesubsection}{1em}{}

\pagestyle{fancy}
\setlength{\headheight}{14pt}
\fancyhf{}
\fancyhead[L]{\small ДЗИ по математика -- Профилирана подготовка}
\fancyhead[R]{\small Справочни материали}
\fancyfoot[C]{\thepage}
\renewcommand{\headrulewidth}{0.4pt}

\setlength{\parindent}{0pt}
\setlength{\parskip}{4pt}

\renewcommand{\tg}{\operatorname{tg}}
\newcommand{\cotg}{\operatorname{cotg}}

\begin{document}

\begin{titlepage}
\centering
\vspace*{1.5cm}

{\large МИНИСТЕРСТВО НА ОБРАЗОВАНИЕТО И НАУКАТА\par}

\vspace{3.2cm}

\begin{tikzpicture}
\draw[titleblue, line width=1pt] (0,0) -- (15.5,0);
\end{tikzpicture}

\vspace{0.8cm}
{\Huge\bfseries\color{titleblue} ДЪРЖАВЕН ЗРЕЛОСТЕН ИЗПИТ\par}
\vspace{0.4cm}
{\Huge\bfseries\color{titleblue} ПО МАТЕМАТИКА\par}

\vspace{1cm}
{\LARGE\itshape ПРОФИЛИРАНА ПОДГОТОВКА\par}

\vspace{0.8cm}
\begin{tikzpicture}
\draw[titleblue, line width=1pt] (0,0) -- (15.5,0);
\end{tikzpicture}

\vspace{1.6cm}
{\Large\bfseries СПРАВОЧНИ МАТЕРИАЛИ\par}

\vfill

\vspace{2cm}

{\normalsize 21.06.2026 г.\par}

\vspace{1cm}

\end{titlepage}

\tableofcontents
\newpage

\section{Вектори и координати}

$$\vec{a} = (a_x, a_y) \qquad \vec{b} = (b_x, b_y)$$

$$|\vec{a}| = \sqrt{a_x^2 + a_y^2}$$

$$\vec{a} \cdot \vec{b} = a_x b_x + a_y b_y$$

$$\vec{a} \cdot \vec{b} = |\vec{a}|\,|\vec{b}|\cos\angle(\vec{a}, \vec{b})$$

\section{Аналитична геометрия в равнината}

$$ax + by + c = 0 \quad \text{— общо уравнение на права}$$

$$y = kx + b \quad \text{— декартово уравнение на права}$$

$$g:\ \frac{x - x_1}{x_2 - x_1} = \frac{y - y_1}{y_2 - y_1},\quad M_1(x_1, y_1),\ M_2(x_2, y_2)$$
\centerline{— уравнение на права през две точки}

$$(x - a)^2 + (y - b)^2 = r^2$$
\centerline{— нормално уравнение на окръжност с център $O(a, b)$ и радиус $r$}

\subsection{Криви от втора степен}

\begin{center}
\begin{tabular}{lll}
\toprule
\textbf{Крива} & \textbf{Канонично уравнение} & \textbf{Фокуси} \\
\midrule
Елипса & $\dfrac{x^2}{a^2} + \dfrac{y^2}{b^2} = 1$ & $F_1(-c,0)$, $F_2(c,0)$, $c^2=a^2-b^2$ \\[10pt]
Хипербола & $\dfrac{x^2}{a^2} - \dfrac{y^2}{b^2} = 1$ & $F_1(-c,0)$, $F_2(c,0)$, $c^2=a^2+b^2$ \\[10pt]
Парабола & $y^2 = 2px$ & $F\left(\dfrac{p}{2},0\right)$ \\[8pt]
\bottomrule
\end{tabular}
\end{center}

\textbf{Ъгъл $\varphi$ между две прави} $g_1: a_1x+b_1y+c_1=0$ и $g_2: a_2x+b_2y+c_2=0$:

$$\cos\varphi = \frac{|a_1a_2+b_1b_2|}{\sqrt{a_1^2+b_1^2}\cdot\sqrt{a_2^2+b_2^2}}$$

\section{Зависимости в триъгълник и успоредник}

\textbf{Правоъгълен триъгълник:}

$$c^2 = a^2+b^2 \qquad S=\frac12 ab=\frac12 c\,h_c$$

$$a^2=a_1c \qquad b^2=b_1c \qquad h_c^2=a_1b_1 \qquad r=\frac{a+b-c}{2}$$

$$\sin\alpha=\frac{a}{c} \qquad \cos\alpha=\frac{b}{c} \qquad \tg\,\alpha=\frac{a}{b} \qquad \cotg\,\alpha=\frac{b}{a}$$

\textbf{Произволен триъгълник:}

$$\frac{a}{\sin\alpha}=\frac{b}{\sin\beta}=\frac{c}{\sin\gamma}=2R$$

$$a^2=b^2+c^2-2bc\cos\alpha \qquad b^2=a^2+c^2-2ac\cos\beta \qquad c^2=a^2+b^2-2ab\cos\gamma$$

\textbf{Формула за медиана:}

$$m_a^2=\tfrac14(2b^2+2c^2-a^2) \qquad m_b^2=\tfrac14(2a^2+2c^2-b^2) \qquad m_c^2=\tfrac14(2a^2+2b^2-c^2)$$

\textbf{Формула за ъглополовяща:}

$$\frac{a}{b}=\frac{n}{m} \qquad l_c^2=ab-mn$$

\textbf{Формула за диагоналите на успоредник:}

$$d_1^2+d_2^2=2a^2+2b^2$$

\section{Формули за лице}

\textbf{Триъгълник:}
$$S=\frac12 c\,h_c \qquad S=\frac12 ab\sin\gamma \qquad S=\sqrt{p(p-a)(p-b)(p-c)} \qquad S=pr \qquad S=\frac{abc}{4R}$$

\textbf{Успоредник:} \quad $S=a\,h_a \qquad S=ab\sin\alpha$

\textbf{Трапец:} \quad $S=\dfrac{a+b}{2}h$

\textbf{Четириъгълник:} \quad $S=\dfrac12 d_1d_2\sin\varphi$, \ $\varphi$ -- ъгъл между диагоналите

\textbf{Описан многоъгълник:} \quad $S=pr$

\section{Ръбести и валчести тела}

\textbf{Права призма:} \quad $S_{\text{л}}=Ph \qquad S=2S_B+S_{\text{л}} \qquad V=Bh$

\textbf{Правилна пирамида:} \quad $S_{\text{л}}=\dfrac{Pa}{2} \qquad S=S_B+S_{\text{л}} \qquad V=\dfrac13 Bh$

\textbf{Пресечена пирамида:}
$$S_{\text{л}}=(P+P_1)\frac{a}{2} \qquad S=S_B+S_{B_1}+S_{\text{л}} \qquad V=\frac{h}{3}\left(B+B_1+\sqrt{B\,B_1}\right)$$

\textbf{Прав кръгов цилиндър:} \quad $S_{\text{л}}=2\pi r h \qquad S=2\pi rh+2\pi r^2 \qquad V=\pi r^2 h$

\textbf{Прав кръгов конус:} \quad $S_{\text{л}}=\pi r l \qquad S=\pi rl+\pi r^2 \qquad V=\dfrac{\pi r^2 h}{3}$

\textbf{Прав кръгов пресечен конус:}
$$S_{\text{л}}=\pi l(R+r) \qquad S=\pi l(R+r)+\pi(R^2+r^2) \qquad V=\frac{\pi h}{3}\left(R^2+Rr+r^2\right)$$

\textbf{Сфера и кълбо:} \quad $S=4\pi r^2 \qquad V=\dfrac43\pi r^3$

\textbf{Лице на проекция} на многоъгълник от равнина $\alpha$ върху равнина $\beta$:
$$S_1=S\cos\varphi$$
където $S$ е лице на многоъгълника от равнината $\alpha$, а $\varphi$ е острият ъгъл между равнините $\alpha$ и $\beta$.

\newpage
\section{Тригонометрични функции}

\begin{center}
\renewcommand{\arraystretch}{1.6}
\begin{tabular}{l*{11}{c}}
\toprule
$\alpha^\circ$ & 0° & 30° & 45° & 60° & 90° & 120° & 135° & 150° & 180° & 270° & 360° \\
\midrule
$\alpha$ rad & 0 & $\frac{\pi}{6}$ & $\frac{\pi}{4}$ & $\frac{\pi}{3}$ & $\frac{\pi}{2}$ & $\frac{2\pi}{3}$ & $\frac{3\pi}{4}$ & $\frac{5\pi}{6}$ & $\pi$ & $\frac{3\pi}{2}$ & $2\pi$ \\
$\sin\alpha$ & 0 & $\frac12$ & $\frac{\sqrt2}{2}$ & $\frac{\sqrt3}{2}$ & 1 & $\frac{\sqrt3}{2}$ & $\frac{\sqrt2}{2}$ & $\frac12$ & 0 & $-1$ & 0 \\
$\cos\alpha$ & 1 & $\frac{\sqrt3}{2}$ & $\frac{\sqrt2}{2}$ & $\frac12$ & 0 & $-\frac12$ & $-\frac{\sqrt2}{2}$ & $-\frac{\sqrt3}{2}$ & $-1$ & 0 & 1 \\
$\tg\,\alpha$ & 0 & $\frac{\sqrt3}{3}$ & 1 & $\sqrt3$ & -- & $-\sqrt3$ & $-1$ & $-\frac{\sqrt3}{3}$ & 0 & -- & 0 \\
$\cotg\,\alpha$ & -- & $\sqrt3$ & 1 & $\frac{\sqrt3}{3}$ & 0 & $-\frac{\sqrt3}{3}$ & $-1$ & $-\sqrt3$ & -- & 0 & -- \\
\bottomrule
\end{tabular}
\end{center}

\subsection{Ъгли}

\begin{center}
\renewcommand{\arraystretch}{1.6}
\begin{tabular}{l*{5}{c}}
\toprule
 & $-\alpha$ & $90^\circ\pm\alpha$ & $180^\circ\pm\alpha$ & $270^\circ\pm\alpha$ & $360^\circ\pm\alpha$ \\
\midrule
$\sin$ & $-\sin\alpha$ & $\cos\alpha$ & $\mp\sin\alpha$ & $-\cos\alpha$ & $\pm\sin\alpha$ \\
$\cos$ & $\cos\alpha$ & $\mp\sin\alpha$ & $-\cos\alpha$ & $\pm\sin\alpha$ & $\cos\alpha$ \\
$\tg$ & $-\tg\,\alpha$ & $\mp\cotg\,\alpha$ & $\pm\tg\,\alpha$ & $\mp\cotg\,\alpha$ & $\pm\tg\,\alpha$ \\
$\cotg$ & $-\cotg\,\alpha$ & $\mp\tg\,\alpha$ & $\pm\cotg\,\alpha$ & $\mp\tg\,\alpha$ & $\pm\cotg\,\alpha$ \\
\bottomrule
\end{tabular}
\end{center}

\subsection{Тригонометрични тъждества}

$$\sin(\alpha\pm\beta)=\sin\alpha\cos\beta\pm\cos\alpha\sin\beta$$
$$\cos(\alpha\pm\beta)=\cos\alpha\cos\beta\mp\sin\alpha\sin\beta$$

$$\sin 2\alpha = 2\sin\alpha\cos\alpha$$
$$\cos 2\alpha = \cos^2\alpha-\sin^2\alpha = 2\cos^2\alpha-1 = 1-2\sin^2\alpha$$

$$\sin\alpha+\sin\beta=2\sin\frac{\alpha+\beta}{2}\cos\frac{\alpha-\beta}{2} \qquad
\sin\alpha-\sin\beta=2\cos\frac{\alpha+\beta}{2}\sin\frac{\alpha-\beta}{2}$$

$$\cos\alpha+\cos\beta=2\cos\frac{\alpha+\beta}{2}\cos\frac{\alpha-\beta}{2} \qquad
\cos\alpha-\cos\beta=-2\sin\frac{\alpha+\beta}{2}\sin\frac{\alpha-\beta}{2}$$

$$1-\cos2\alpha=2\sin^2\alpha \qquad 1+\cos2\alpha=2\cos^2\alpha$$

$$\sin2\alpha=\frac{2\tg\,\alpha}{1+\tg^2\alpha} \qquad \cos2\alpha=\frac{1-\tg^2\alpha}{1+\tg^2\alpha}$$

\section{Полиноми на една променлива}

\textbf{Теорема на Безу:}
$$P_n(x)=(x-x_0)Q_{n-1}(x)+P_n(x_0)$$

\textbf{Схема на Хорнер} за делене на $P_n(x)=a_0x^n+a_1x^{n-1}+\cdots+a_n$ на $(x-x_0)$:

\begin{center}
\renewcommand{\arraystretch}{1.6}
\begin{tabular}{l*{6}{c}}
\toprule
 & $a_0$ & $a_1$ & $a_2$ & $\ldots$ & $a_{n-1}$ & $a_n$ \\
\midrule
$x_0$ & $b_0{=}a_0$ & $b_1{=}x_0b_0{+}a_1$ & $b_2{=}x_0b_1{+}a_2$ & $\ldots$ & $b_{n-1}{=}x_0b_{n-2}{+}a_{n-1}$ & $x_0b_{n-1}{+}a_n{=}P_n(x_0)$ \\
\bottomrule
\end{tabular}
\end{center}

\section{Числови редици. Граници на редици}

\textbf{Нютонов бином:}
$$(a+b)^n=\binom{n}{0}a^n+\binom{n}{1}a^{n-1}b+\cdots+\binom{n}{k}a^{n-k}b^k+\cdots+\binom{n}{n}b^n$$

$$\binom{n}{k}=C_n^k \qquad \binom{n}{k}=\binom{n}{n-k} \qquad \binom{n}{k}=\binom{n-1}{k-1}+\binom{n-1}{k}$$

\textbf{Сума на безкрайно намаляваща геометрична прогресия:}
$$S=\frac{a_1}{1-q},\quad |q|<1$$

\textbf{Граници на редици:}
$$\lim_{n\to\infty}\frac1n=0 \qquad \lim_{n\to\infty}q^n=0,\ |q|<1 \qquad \lim_{n\to\infty}\left(1+\frac1n\right)^n=e$$

\section{Функции. Непрекъснатост и диференцируемост}

\textbf{Граници на функции:}
$$\lim_{x\to\infty}\frac1x=0 \qquad \lim_{x\to0}\frac{\sin x}{x}=1$$
$$\lim_{x\to\infty}\left(1+\frac1x\right)^x=e \qquad \lim_{x\to0}(1+x)^{\frac1x}=e$$

Ако $\displaystyle\lim_{x\to c}f(x)=0$, то $\displaystyle\lim_{x\to c}\frac{\sin f(x)}{f(x)}=1$.

$$\lim_{x\to\infty}\frac{a_0x^n+a_1x^{n-1}+\cdots+a_n}{b_0x^m+b_1x^{m-1}+\cdots+b_{m-1}}=
\begin{cases}
\infty & \text{при } n>m \\[4pt]
\dfrac{a_0}{b_0} & \text{при } n=m \\[4pt]
0 & \text{при } n<m
\end{cases}$$

Функцията $f(x)$ е непрекъсната в точката $x_0$ от дефиниционното ѝ множество, ако
$\displaystyle\lim_{x\to x_0}f(x)=f(x_0)$.

\subsection{Производни на някои функции}

\begin{center}
\renewcommand{\arraystretch}{1.7}
\begin{tabular}{ll}
\toprule
\textbf{Функция} & \textbf{Производна} \\
\midrule
$c$ (константа) & $0$ \\
$x$ & $1$ \\
$x^a$ & $a\,x^{a-1}$ \\
$\sin x$ & $\cos x$ \\
$\cos x$ & $-\sin x$ \\
$\tg\,x$ & $\dfrac{1}{\cos^2 x}$ \\
$\cotg\,x$ & $-\dfrac{1}{\sin^2 x}$ \\
$e^x$ & $e^x$ \\
$a^x\ (a>0)$ & $a^x\ln a$ \\
$\ln x$ & $\dfrac1x$ \\
$\log_a x\ (a{>}0,a{\neq}1,x{>}0)$ & $\dfrac{1}{x\ln a}$ \\
\bottomrule
\end{tabular}
\end{center}

\textbf{Правила за диференциране:}

$$(cf(x))'=c\,f'(x) \qquad (f(x)\pm g(x))'=f'(x)\pm g'(x)$$
$$(f(x)g(x))'=f'(x)g(x)+f(x)g'(x)$$
$$\left(\frac{f(x)}{g(x)}\right)'=\frac{f'(x)g(x)-f(x)g'(x)}{g^2(x)}$$
$$\big(g(f(x))\big)'=g'(f(x))\cdot f'(x)$$

\section{Приложения на математическия анализ}

\textbf{Допирателна} към графиката на функция $y=f(x)$ в точката $(x_0,f(x_0))$:
$$y-f(x_0)=f'(x_0)(x-x_0)$$

\textbf{Допирателна $t$ към окръжност} $x^2+y^2=R^2$ в точка $M(x_0,y_0)$ от нея:
$$t:\ x_0x+y_0y=R^2$$

\textbf{Допирателна $t$ към крива от втора степен} в точка $M(x_0,y_0)$ от нея:

Елипса: \quad $t:\ \dfrac{x_0x}{a^2}+\dfrac{y_0y}{b^2}=1$ \qquad\qquad Хипербола: \quad $t:\ \dfrac{x_0x}{a^2}-\dfrac{y_0y}{b^2}=1$

\section{Комбинаторика}

\begin{center}
\small
\renewcommand{\arraystretch}{2.2}
\begin{tabularx}{\textwidth}{X>{\centering\arraybackslash}X>{\centering\arraybackslash}X}
\toprule
 & \textbf{без повторение от $n$ елемента от $k$-ти клас} & \textbf{с повторение от $n$ елемента от $k$-ти клас} \\
\midrule
Пермутации & $P_n=n!$ & $P_n(n_1,\ldots,n_k)=\dfrac{n!}{n_1!\cdots n_k!}$, \newline $\sum n_i=n$ \\
Вариации & $V_n^k=\dfrac{n!}{(n-k)!}$ & $\overline{V}_n^k=n^k$ \\
Комбинации & $C_n^k=\dfrac{n!}{k!(n-k)!}$ & $\overline{C}_n^k=\dfrac{(n+k-1)!}{k!(n-1)!}$ \\
\bottomrule
\end{tabularx}
\end{center}

\section{Вероятности}

\textbf{Класическа вероятност:}
$$P(A)=\frac{\text{брой на благоприятните изходи}}{\text{общ брой на изходите}},\quad 0\le P(A)\le1,\quad P(\emptyset)=0,\ P(\Omega)=1$$

\textbf{Вероятност на сума:}
$$P(A\cup B)=P(A)+P(B),\quad A\text{ и }B\text{ са несъвместими}$$
$$P(A\cup B)=P(A)+P(B)-P(AB),\quad A\text{ и }B\text{ са съвместими}$$

\textbf{Условна вероятност:}
$$P(A\mid B)=\frac{P(AB)}{P(B)},\quad P(B)\neq0$$

\textbf{Формула за пълната вероятност:}

Нека $B_1,B_2,\ldots,B_n$ е пълна група събития при даден опит. Тогава вероятността да настъпи случайното събитие $A$ е:
$$P(A)=P(A\mid B_1)P(B_1)+P(A\mid B_2)P(B_2)+\cdots+P(A\mid B_n)P(B_n)$$

\textbf{Формула на Бейс:} за всяко $k=1,2,\ldots,n$ е изпълнено
$$P(B_k\mid A)=\frac{P(A\mid B_k)\,P(B_k)}{P(A)}$$

\section{Математическо очакване, дисперсия и стандартно отклонение}

\begin{center}
\begin{tabular}{l*{5}{c}}
\toprule
$X$ & $x_1$ & $x_2$ & $\ldots$ & $x_k$ \\
$P$ & $p_1$ & $p_2$ & $\ldots$ & $p_k$ \\
\bottomrule
\end{tabular}
\end{center}

\textbf{Математическо очакване:}
$$E(X)=x_1p_1+x_2p_2+\cdots+x_kp_k$$

\textbf{Дисперсия:}
$$D(X)=(x_1-E(X))^2p_1+(x_2-E(X))^2p_2+\cdots+(x_k-E(X))^2p_k$$
$$D(X)=E(X^2)-(E(X))^2=E\big[(X-EX)^2\big]$$

\textbf{Стандартно отклонение:} \quad $\sigma=\sqrt{D(X)}$

\textbf{Биномно разпределение} с параметри $n$, $p$ и $q=1-p$:

\begin{center}
\renewcommand{\arraystretch}{1.6}
\begin{tabular}{l*{6}{c}}
\toprule
$X$ & 0 & 1 & $\ldots$ & $k$ & $\ldots$ & $n$ \\
\midrule
$P$ & $C_n^0p^0q^n$ & $C_n^1p^1q^{n-1}$ & $\ldots$ & $C_n^kp^kq^{n-k}$ & $\ldots$ & $C_n^np^nq^0$ \\
\bottomrule
\end{tabular}
\end{center}

$$E(X)=np \qquad D(X)=npq \qquad \sigma=\sqrt{npq}$$

\section{Нормално разпределение}

Нормално разпределение $N(\mu,\sigma^2)$ на случайна величина $X$.

\textbf{Функция на плътност:}
$$f(x)=\frac{1}{\sigma\sqrt{2\pi}}\,e^{-\frac12\left(\frac{x-\mu}{\sigma}\right)^2},\quad x\in(-\infty,\infty)$$

където $E(X)=\mu$ -- математическо очакване, $\sigma$ -- стандартно отклонение.

\textbf{Стандартно нормално разпределение} $N(0,1)$ на случайна величина $Z$:
$$Z=\frac{X-\mu}{\sigma}$$

\subsection{Таблица за стойностите на стандартното нормално разпределение $N(0,1)$}

Стойности на площите под нормалната крива, вляво от съответната стойност на аргумента:
$$P(Z<z)=\Phi(z),\quad \text{когато } z\ge0$$
За отрицателни стойности на $z$ площите се намират чрез $1-\Phi(z)$.

\begin{center}
\small
\renewcommand{\arraystretch}{1.2}
\begin{tabular}{c*{10}{c}}
\toprule
$z$ & 0,00 & 0,01 & 0,02 & 0,03 & 0,04 & 0,05 & 0,06 & 0,07 & 0,08 & 0,09 \\
\midrule
0,0 & 0,5000 & 0,5040 & 0,5080 & 0,5120 & 0,5160 & 0,5199 & 0,5239 & 0,5279 & 0,5319 & 0,5359 \\
0,1 & 0,5398 & 0,5438 & 0,5478 & 0,5517 & 0,5557 & 0,5596 & 0,5636 & 0,5675 & 0,5714 & 0,5753 \\
0,2 & 0,5793 & 0,5832 & 0,5871 & 0,5910 & 0,5948 & 0,5987 & 0,6026 & 0,6064 & 0,6103 & 0,6141 \\
0,3 & 0,6179 & 0,6217 & 0,6255 & 0,6293 & 0,6331 & 0,6368 & 0,6406 & 0,6443 & 0,6480 & 0,6517 \\
0,4 & 0,6554 & 0,6591 & 0,6628 & 0,6664 & 0,6700 & 0,6736 & 0,6772 & 0,6808 & 0,6844 & 0,6879 \\
0,5 & 0,6915 & 0,6950 & 0,6985 & 0,7019 & 0,7054 & 0,7088 & 0,7123 & 0,7157 & 0,7190 & 0,7224 \\
0,6 & 0,7257 & 0,7291 & 0,7324 & 0,7357 & 0,7389 & 0,7422 & 0,7454 & 0,7486 & 0,7517 & 0,7549 \\
0,7 & 0,7580 & 0,7611 & 0,7642 & 0,7673 & 0,7704 & 0,7734 & 0,7764 & 0,7794 & 0,7823 & 0,7852 \\
0,8 & 0,7881 & 0,7910 & 0,7939 & 0,7967 & 0,7995 & 0,8023 & 0,8051 & 0,8078 & 0,8106 & 0,8133 \\
0,9 & 0,8159 & 0,8186 & 0,8212 & 0,8238 & 0,8264 & 0,8289 & 0,8315 & 0,8340 & 0,8365 & 0,8389 \\
1,0 & 0,8413 & 0,8438 & 0,8461 & 0,8485 & 0,8508 & 0,8531 & 0,8554 & 0,8577 & 0,8599 & 0,8621 \\
1,1 & 0,8643 & 0,8665 & 0,8686 & 0,8708 & 0,8729 & 0,8749 & 0,8770 & 0,8790 & 0,8810 & 0,8830 \\
1,2 & 0,8849 & 0,8869 & 0,8888 & 0,8907 & 0,8925 & 0,8944 & 0,8962 & 0,8980 & 0,8997 & 0,9015 \\
1,3 & 0,9032 & 0,9049 & 0,9066 & 0,9082 & 0,9099 & 0,9115 & 0,9131 & 0,9147 & 0,9162 & 0,9177 \\
1,4 & 0,9192 & 0,9207 & 0,9222 & 0,9236 & 0,9251 & 0,9265 & 0,9279 & 0,9292 & 0,9306 & 0,9319 \\
1,5 & 0,9332 & 0,9345 & 0,9357 & 0,9370 & 0,9382 & 0,9394 & 0,9406 & 0,9418 & 0,9429 & 0,9441 \\
1,6 & 0,9452 & 0,9463 & 0,9474 & 0,9484 & 0,9495 & 0,9505 & 0,9515 & 0,9525 & 0,9535 & 0,9545 \\
1,7 & 0,9554 & 0,9564 & 0,9573 & 0,9582 & 0,9591 & 0,9599 & 0,9608 & 0,9616 & 0,9625 & 0,9633 \\
1,8 & 0,9641 & 0,9649 & 0,9656 & 0,9664 & 0,9671 & 0,9678 & 0,9686 & 0,9693 & 0,9699 & 0,9706 \\
1,9 & 0,9713 & 0,9719 & 0,9726 & 0,9732 & 0,9738 & 0,9744 & 0,9750 & 0,9756 & 0,9761 & 0,9767 \\
2,0 & 0,9772 & 0,9778 & 0,9783 & 0,9788 & 0,9793 & 0,9798 & 0,9803 & 0,9808 & 0,9812 & 0,9817 \\
2,1 & 0,9821 & 0,9826 & 0,9830 & 0,9834 & 0,9838 & 0,9842 & 0,9846 & 0,9850 & 0,9854 & 0,9857 \\
2,2 & 0,9861 & 0,9864 & 0,9868 & 0,9871 & 0,9875 & 0,9878 & 0,9881 & 0,9884 & 0,9887 & 0,9890 \\
2,3 & 0,9893 & 0,9896 & 0,9898 & 0,9901 & 0,9904 & 0,9906 & 0,9909 & 0,9911 & 0,9913 & 0,9916 \\
2,4 & 0,9918 & 0,9920 & 0,9922 & 0,9925 & 0,9927 & 0,9929 & 0,9931 & 0,9932 & 0,9934 & 0,9936 \\
2,5 & 0,9938 & 0,9940 & 0,9941 & 0,9943 & 0,9945 & 0,9946 & 0,9948 & 0,9949 & 0,9951 & 0,9952 \\
2,6 & 0,9953 & 0,9955 & 0,9956 & 0,9957 & 0,9959 & 0,9960 & 0,9961 & 0,9962 & 0,9963 & 0,9964 \\
2,7 & 0,9965 & 0,9966 & 0,9967 & 0,9968 & 0,9969 & 0,9970 & 0,9971 & 0,9972 & 0,9973 & 0,9974 \\
2,8 & 0,9974 & 0,9975 & 0,9976 & 0,9977 & 0,9977 & 0,9978 & 0,9979 & 0,9979 & 0,9980 & 0,9981 \\
2,9 & 0,9981 & 0,9982 & 0,9982 & 0,9983 & 0,9984 & 0,9984 & 0,9985 & 0,9985 & 0,9986 & 0,9986 \\
3,0 & 0,9987 & 0,9987 & 0,9987 & 0,9988 & 0,9988 & 0,9989 & 0,9989 & 0,9989 & 0,9990 & 0,9990 \\
\bottomrule
\end{tabular}
\end{center}

\end{document}
